\paragraph{R3-W1 (Techniques reduce to a common pipeline; apparent contradiction; prototype uses third-party libraries such as FAISS).}

\begin{itemize}[leftmargin=1.2em, itemsep=1pt, topsep=2pt]
\item \textit{W1: ``The proposed techniques \dots reduce to a common pipeline: embedding, index construction, nearest-neighbour search, and utility-based subset selection \dots the authors simultaneously discuss that these elements can be handled as separate stages within existing systems. This contradiction weakens the case \dots the experiment in DcCA also treated the criticized stages independently with third-party libraries, e.g., using FAISS \dots''}
\end{itemize}
\smallskip

\jp{This response argues against the reviewer instead of helping the reviewer understand the paper. The response is hard to read/understand.}

\revision{We appreciate this insightful observation. The apparent contradiction arises because our paper has two different goals: (1) to articulate a \emph{vision} for future data management systems, and (2) to build a proof-of-concept prototype using today's available systems. There is therefore no contradiction between our vision and the current implementation.

\textbf{The vision.} The central thesis of the paper is that counterfactual querying should become a \emph{first-class data management primitive}. Existing systems expose vector search, filtering, reranking, and subset selection as independent operators. Our argument is that the opposite-outcome predicate, feasibility constraints, intervention cost $\omega$, and set objective $\Omega$ should instead become part of the declarative query semantics, so that a database optimizer can reason about them jointly during query optimization, index construction, and execution planning.

\textbf{The prototype.} Current database systems do not yet support such an abstraction. Consequently, our prototype is intentionally implemented using existing components, such as FAISS for candidate retrieval, together with feasibility validation and set optimization. FAISS is used only as a proposal mechanism for generating candidate counterfactuals; it does not implement the semantics advocated in this paper. The remaining components approximate the desired query semantics outside the retrieval engine because current systems provide no mechanism to optimize them jointly.

Our criticism is therefore not directed at using FAISS or other mature retrieval components. Rather, it is directed at the limitation of current database systems, which treat these components as separate stages. The prototype demonstrates that counterfactual querying is practical with today's infrastructure, while the vision identifies the next research challenge: elevating counterfactual querying from an application-level workflow to a native database abstraction in which feasibility, intervention cost, and set-level objectives are optimized together rather than handled by independent stages.

To avoid this confusion, we revised Section~5.1 to explicitly distinguish the current proof-of-concept implementation from the long-term vision. We also clarify that FAISS serves only as the proposal mechanism, whereas the research challenge is to develop future query processors and indexes that natively support the semantics of counterfactual querying.}

\noindent\textbf{Response.}
We resolve the apparent contradiction by stating the thesis precisely: existing systems \emph{can only} treat feasibility, cost, and set objectives as separate stages, and that staging is exactly what is insufficient. As Section~5.1 argues, \emph{``feasibility constraints and opposite-outcome predicates reshape the search space rather than act as simple filters, while cost-faithful metrics (P2) often deviate from geometric distances assumed by existing indexes,''} so systems \emph{``may return infeasible candidates or rely on post hoc validation, weakening both efficiency (P3) and correctness.''} The contribution is to make $\omega$, the counterfactual predicate, and $\Omega$ first-class and jointly optimized---not a new ANN index. FAISS in the prototype implements only the \emph{proposal} step (Step~2.a) under a fixed metric; the paper itself states that similarity search \emph{``serves only as a proposal mechanism, while correctness is determined through \dots\ validation.''} Using mature components for the proposal layer is honest scaffolding for a proof-of-concept. Our new quality study (Table~\ref{tab:quality}) shows why the ``utility-based subset selection'' stage is not interchangeable with generic top-$k$: top-$k$ is optimal for none of the set-level objectives.

\noindent\emph{Revision to Section~5.1:}
\begin{quote}
\textcolor{blue}{We emphasize that decomposing counterfactual querying into independent stages---the only option in current systems---is precisely the limitation we identify, not an endorsement: post-hoc feasibility filtering reshapes the feasible space too late and degrades recall and stability, the embedding metric is not the cost-faithful $\omega$, and tuple-at-a-time selection cannot optimize the coupled objective $\Omega$. Our prototype uses FAISS only as a proposal mechanism (Step~2.a); folding $\omega$ and $\Omega$ into retrieval and selection is the research agenda this vision defines.}
\end{quote}
We also add the positioning table (Table~\ref{tab:positioning}), which contrasts DcCA with constrained retrieval, vector search, nearest-neighbor search, subset selection, and the feasibility-aware FACE~\cite{poyiadzi2020face}.
